\section{Scope: what is and is not proved about $A(p)$}
\label{a:sec:scope}
Before any claim about closed forms is stated, the phrase itself must be made precise: at least four inequivalent readings are in play, and a result for one is not a result for another.
\begin{enumerate}[label=\textup{(\Alph*)},leftmargin=*]
\item[\textbf{(V)}] the \emph{value} $A(p_0)$ is elementary or algebraic, for a fixed rational $p_0$;
\item[\textbf{(E)}] the \emph{map} $p\mapsto A(p)$ is an elementary function of $p$;
\item[\textbf{(D)}] the map is \emph{differentially algebraic in $p$}, i.e.\ satisfies some algebraic ODE in $p$;
\item[\textbf{(S)}] the map is expressible via standard special functions ($\Gamma$, $\zeta$, ${}_2F_1$, theta or $q$-series).
\end{enumerate}
\Cref{a:thm:main-ht} speaks to the function $z\mapsto\psi_r(z)$ at fixed $r$, and licenses exactly this much: \emph{no evaluation of $A(p)$ can be short-circuited through an algebraic differential equation for the lineariser, at any subcritical parameter.} It does not, by itself, decide any of (V), (E), (D) or (S). Two gaps must be stated plainly.

\paragraph{The value gap.}
A hypertranscendental function can take elementary --- even algebraic --- values at individual points: $\Gamma$ is hypertranscendental by H\"older's theorem, yet $\Gamma(\tfrac12)=\sqrt{\pi}$. Nothing so far forbids, say, $A(\tfrac14)$ from being an algebraic number. There \emph{is} a genuine values theorem in this circle of ideas: Becker and Bergweiler proved in 1993 \cite{Becker1993} that the B\"ottcher conjugacy between two polynomials of the \emph{same degree} $d\ge2$ with algebraic coefficients (not linearly conjugate to monomials or Chebyshev polynomials) is transcendental and takes transcendental values at algebraic points; Mahler's method drives the proof (for the method, see Nishioka \cite{Nishioka1996}). One might hope to transfer this to the Koenigs value $\psi_r(\tfrac12)$. It does not transfer, for two independent reasons. \emph{Formally}, the Schr\"oder equation pairs an inner map of degree two with an outer map --- multiplication by $r$ --- of degree one, so the equal-degree hypothesis fails outright. \emph{Structurally}, no easy repair exists: Mahler's method balances doubly-exponential analytic decay along an orbit against comparable height growth, and in the B\"ottcher setting both exponents run like $d^{\,n}$; in the attracting Koenigs setting the orbit $w_n=f_r^{\circ n}(\tfrac12)$ decays only geometrically, $|w_n|\sim A\,r^{n}$, while the heights of the rational numbers $w_n$ still grow doubly exponentially, $h(w_n)\sim C\,2^{\,n}$, and the resulting inequalities run the wrong way at every $n$. Note that the obstruction is \emph{not} the germ/basin issue --- the pullback of \Cref{a:rem:basin} moves the evaluation point arbitrarily deep into the germ at the cost of a rational factor --- it is the degree structure of the equation itself. The consequence deserves emphasis:
\begin{center}
\emph{transcendence --- indeed, even irrationality --- of $A(p_0)$ is at present not proved\\ for a single rational $p_0\in(0,\tfrac12)$.}
\end{center}

\paragraph{The parameter gap.}
\Cref{a:thm:main-ht} says nothing about the map $p\mapsto A(p)$: readings (E) and (D) are open. To the best of our knowledge the question is not merely unsolved but unaddressed in the literature. The natural machinery --- the parametrised differential Galois theory of difference equations of Hardouin and Singer \cite{HardouinSinger2008}, designed precisely to detect differential relations in auxiliary directions for linear difference equations such as $\sigma\psi=r\psi$ --- does not apply as it stands, because the derivation $\partial_r$ fails to commute with the substitution operator $\sigma\colon w\mapsto f_r(w)$, which itself moves with $r$. Nor does any feature of the near-critical asymptotics obstruct (E) or (D): logarithmic corrections to \eqref{a:eq:Aasym} are themselves elementary functions of $\varepsilon$.

\subsection{Inverse-symbolic evidence at eleven rational parameters}
\label{a:subsec:pslq}

Since the theorems leave (V) open, the value question was attacked directly by integer-relation detection. The values
\[
A(p),\qquad p\in\bigl\{\tfrac1{10},\tfrac18,\tfrac16,\tfrac15,\tfrac14,\tfrac3{10},\tfrac13,\tfrac38,\tfrac25,\tfrac5{12},\tfrac9{20}\bigr\},
\]
were computed to more than $1150$ significant digits by two mechanically independent routes --- the exact product identity \eqref{a:eq:prodFormula} iterated to convergence, and a pullback of the orbit into the Koenigs germ followed by evaluation of the Koenigs series --- which agreed to at least $1192$ digits at every parameter. Each value was then subjected to a PSLQ battery at $\ge200$-digit effective rigour:
\begin{itemize}
\item \emph{algebraicity}: no integer polynomial relation of degree $\le10$ and coefficient height $\le10^{6}$;
\item \emph{linearity}: no rational linear relation, at height $\le10^{8}$, against the constants $\pi$, $\pi^{2}$, $\ln2$, $\ln3$, $\ln5$, $\gamma$, $\zeta(3)$, $\sqrt2$, $\sqrt3$, $\sqrt5$, $e$, $e^{\pi}$, $\Gamma(\tfrac13)$, $\Gamma(\tfrac14)$ and Catalan's constant (also tested for $1/A$, the quasi-stationary mean itself, and for $A/\varepsilon$);
\item \emph{bilinearity}: no representation $A=L_1/L_2$ with $L_1,L_2$ integer linear forms of height $\le10^{8}$ in a twelve-constant basis (run at $460$ digits);
\item \emph{multiplicativity}: no relation $A=e^{q_0}2^{q_1}3^{q_2}5^{q_3}7^{q_4}\pi^{q_5}$ with rational exponents of height $\le10^{10}$;
\item \emph{cross-parameter}: no rational linear or multiplicative relation linking the eleven values to one another (run at $500$ digits).
\end{itemize}
In all, $123$ tests returned no relation --- indeed no discovery-stage candidate: every PSLQ run terminated on its internal norm bound, so each null carries a genuine ``no relation up to the stated height'' guarantee. Throughout, the precision head-room (digits divided by basis size) exceeded every height cap by at least eight orders of magnitude, so a returned relation would have been meaningful rather than numerical noise.\footnote{Computations in \texttt{mpmath} at working precisions of $310$--$1200$ digits; scripts, the $1150$-digit values, and the full per-test log are archived with the project sources.}

This is finite evidence, and its scope must be stated honestly: it excludes low-complexity closed forms over the tested constants, and nothing more. A closed form over untested constants ($\Gamma(\tfrac17)$, say, or periods of higher-genus curves), or one of height above the caps, would have escaped the search. It is evidence for the conjecture, not a proof.

\subsection{The working claim}

\begin{quote}
\textit{For every fixed $r=2p\in(0,1)$, the Koenigs linearisation $\psi_r$ is hypertranscendental in its own variable (\Cref{a:thm:main-ht}) --- a theorem, with no exceptional parameters in the attracting window, since differentially algebraic Schr\"oder functions exist only over repelling fixed points. Consequently no elementary expression in $z$ for the conjugacy can be evaluated to produce $A(p)$. Whether individual values $A(p_0)$ are transcendental (or even irrational), and whether the map $p\mapsto A(p)$ is elementary or differentially algebraic in $p$, remain open questions; no closed form under any reading is known, none is expected, and eleven rational values resist inverse-symbolic identification against standard constants at over $200$-digit rigour.}
\end{quote}

Full derivations of the elementary solutions at $r=2$ and $r=4$, and their place inside the Becker--Bergweiler classification, are collected in \cref{a:app:integrable}.
